\section{Tipo de Pesquisa}
É chamada pesquisa aplicada aquela que gera conhecimento visando aplicação prática, direcionada para a solução de problemas específicos \cite{sienaMetodologiaPesquisaCientifica2024}. Assim, o presente trabalho busca investigar o desempenho, no que diz respeito à atenção e engajamento dos alunos submentidos a aulas expositivas em ambiente remoto; aplicando e validando o sistema de informação ESPEON, para que,- por meio de relatórios quantitativos - seja possível fundamentar argumentos contrários ao uso da metodologia expositiva, sobretudo em sala de aula remota.

\section{Objetivos da Pesquisa}
Trata-se de uma pesquisa descritiva explicativa, que busca, por meio da observação sistemática, descrever comportamentalmente a população (alunos) submetida à metodologia expositiva em ambiente remoto. Pretende-se, com o auxílio do sistema ESPEON para coleta e análise de dados, evidenciar o baixo desempenho dos alunos expostos a esse tipo de abordagem, de forma a corroborar a presente tese.

É descritiva porque objetiva, primeiramente, caracterizar uma população submetida a um fenômeno (aula expositiva remota), utilizando técnicas padronizadas de coleta de dados (\textit{payload} da API), conforme prescrito por \cite{sienaMetodologiaPesquisaCientifica2024}.

Também é considerada explicativa porque busca, por meio das evidências (relatórios de aula), indicar os principais comportamentos demonstrados pelos alunos, relacionando-os à falta de atenção e engajamento nesse tipo de metodologia. Nesse sentido, reforça-se que:

\begin{quote}
    ``A pesquisa explicativa visa explicar a razão dos fatos, por meio da identificação e
    análise das relações de causa e efeito dos fenômenos. Nas ciências sociais, para a coleta de
    dados pode ser utilizado o método observacional ou outros métodos aplicáveis, a depender do
    desenho da pesquisa, dos objetivos e das hipóteses ou suposições formuladas pelo pesquisador.''
    \cite{sienaMetodologiaPesquisaCientifica2024}
\end{quote}

\section{Abordagem da Pesquisa}
No que se refere à abordagem do problema, Ribeiro et al. (2015) \cite{ribeiroManagementAccountingAnalysis2015} destaca que uma pesquisa quantitativa é aquela que emprega a quantificação tanto na coleta quanto no tratamento dos dados, por meio de técnicas estatísticas. Neste trabalho, serão utilizadas ferramentas de estatística descritiva, como a média aritmética (medida de tendência central), o desvio-padrão (indicador de dispersão) e a moda (valor de maior frequência), de forma a caracterizar a população estudada. Os cálculos serão processados pelo ESPEON no \textit{backend} e consolidados em gráficos, visando facilitar a interpretação dos resultados e possibilitar análises mais claras.

A escolha da estatística descritiva justifica-se pelo fato de que ela permite sintetizar e interpretar grandes volumes de dados de maneira objetiva, fornecendo subsídios para a comparação entre diferentes comportamentos observados durante as aulas remotas.

% Talvez caiba colocar uma definição desses caras na literatura, e um parágrafo final justificando sua escolha

\section{Método de Pesquisa}
Buscando averiguar a problemática do emprego da metodologia expositiva no ambiente remoto de ensino, parte-se da hipótese de que alunos submetidos a esse tipo de aula, seja no EAD ou no ERE, tendem a apresentar dispersão de foco e baixo engajamento.

Para testar essa hipótese, a pesquisa será conduzida por meio de experimentos controlados com alunos do curso de Sistemas de Informação do CEFET/RJ, campus Maria da Graça. Os participantes serão submetidos às condições mencionadas e terão seu comportamento monitorado pelo sistema ESPEON, mediante consentimento prévio, de forma a possibilitar a coleta de dados para posterior análise estatística.

O método adotado é caracterizado como hipotético-dedutivo, visto que, conforme Ribeiro et al. (2015) \cite{ribeiroManagementAccountingAnalysis2015}, fundamenta-se na experimentação para verificar a validade de uma teoria — neste caso, a de que aulas expositivas em ambientes remotos se mostram ineficazes.

\section{Objeto de Estudo}
O objeto de estudo desta pesquisa é o emprego de um sistema de informação para o monitoramento de estudantes submetidos à metodologia expositiva em ambientes remotos de ensino. As contribuições esperadas incluem a identificação de gargalos nesse cenário, possibilitando ao docente adequar sua metodologia à realidade dos alunos. Além disso, o estudo suscita uma reflexão relevante: até que ponto faz sentido utilizar um método tradicional em um contexto fortemente marcado pelo avanço tecnológico?

\section{Procedimentos e Instrumentos de Coleta de Dados}

% \section{Análise de Dados}
